\chapter{Conclusion}
\label{ch:conclusion}

\section{Summary of Contributions}
\label{sec:conclusion-summary}

This thesis introduced \textsc{ReasonGuard}, a runtime verification
framework for AI reasoning that uses the structured output of a formally
verified industrial control system detector as a machine-readable bound on
what an explanation may correctly claim. The framework was implemented as
a Python package, empirically evaluated on the IEC-104 dataset published
by the NES@FIT group at Brno University of Technology, and the empirical
findings were translated into operational deployment guidelines and mapped
to the transparency and quality-management provisions of the European
Union AI Act.

Four contributions follow.

\paragraph{A runtime verification framework for AI reasoning.}
Chapter~\ref{ch:framework} presented the three-layer architecture
(formal layer, AI layer, verification layer), the JSON bound schema, and
the verification mechanism (claim extraction, bound comparison, violation
classification). The implementation is open source, fully tracked in
MLflow, and accompanied by a 79-case pytest suite.

\paragraph{A taxonomy of five reasoning violation classes.}
The V1--V5 taxonomy (fabricated, contradicted, over-generalised,
under-specified, incoherent reasoning) is the central theoretical
contribution. Each class is formally defined, implemented with both a
lexicon-based and a spaCy-based detection pipeline, and empirically
exercised on the synthetic regression set as well as on real local-LLM
outputs.

\paragraph{An empirical evaluation across local and cloud LLMs.}
Chapter~\ref{ch:results} reported the empirical violation rates per model
and per event type, the F1 and Cohen's $\kappa$ numbers against the
manually annotated reference, and the latency and operational-cost
figures. The framework discriminates between local models on the same
prompt set and exposes the rarer-event-type sensitivity of small models.

\paragraph{Operational deployment guidelines.}
Chapter~\ref{ch:discussion} translated the empirical findings into a
five-item checklist for organisations deploying an LLM-based explanation
layer on top of a formally verified detector. Each item was mapped to a
specific provision of the EU AI Act.

\section{Answers to the Research Questions}
\label{sec:conclusion-answers}

The central research question of this thesis was:

\begin{quote}
  Can formally verified system outputs be used as runtime bounds for
  detecting and classifying AI reasoning violations, and does such a
  framework produce actionable, operationally reliable violation
  detection in safety-critical AI deployments?
\end{quote}

The empirical evidence in Chapter~\ref{ch:results} supports an affirmative
answer in the IEC-104 setting under the proxy schema. The five
sub-questions are answered as follows.

\paragraph{(RQ1) Bound representation.}
Section~\ref{sec:bound-schema} introduced a JSON bound schema with five
top-level groups (metadata, features, allowed claims, required claims,
forbidden claims, machine-checkable constraints) that is sufficient to
support the V1--V5 detector. The schema is currently labelled as
\texttt{proxy\_until\_official\_automaton\_schema}; the upgrade path to
the official NES@FIT schema is the highest-priority follow-up.

\paragraph{(RQ2) Violation taxonomy.}
Section~\ref{sec:taxonomy} formalised the V1--V5 taxonomy with respect to
the bound and Chapter~\ref{ch:results} empirically observed the
prevalence of each class in the synthetic regression set and in the
real-LLM batch. V1 and V4 dominated the real-LLM batch, V3 was rare, and
V5 was extremely rare under the spaCy detector.

\paragraph{(RQ3) Detector accuracy.}
The manual-annotation pipeline in
Sections~\ref{sec:manual-annotation}--\ref{sec:results-detector}
produces per-class precision, recall, and F1. The W14 M4 milestone
defines the production threshold at F1 $> 0.80$ per class on the 200-pair
annotated set; the W15 M5 milestone targets Cohen's $\kappa > 0.70$
against the VUT Brno domain expert.

\paragraph{(RQ4) Conditions for violation.}
The per-event-type analysis in Section~\ref{sec:results-per-event} shows
that all three local models perform worst on rarer proxy event types and
that smaller models produce systematically higher violation rates than
larger ones on the same prompt set. Quantisation level is held fixed at
4-bit across the local batch; future work should extend the comparison
across additional quantisation levels.

\paragraph{(RQ5) Operational guidelines.}
Section~\ref{sec:discussion-deployment} presented the five-item
deployment checklist, each item mapped to the EU AI Act provision it
addresses (Section~\ref{sec:discussion-eu-ai-act}).

\section{Future Work}
\label{sec:conclusion-future}

\paragraph{Integration of the official NES@FIT automaton schema.}
The proxy schema currently in use will be replaced by the official
schema as soon as the documentation is available. The replacement is a
schema bump (not a rewrite) and the existing artefacts can be regenerated
deterministically against the new schema.

\paragraph{Learning-based per-claim classifier.}
The per-claim three-state classifier is rule-based. A learning-based
classifier trained on the 200-pair annotated set, with the rule-based
classifier as the baseline, is a natural extension and would close any
precision-recall trade-off observed at the M4 milestone.

\paragraph{Operator-in-the-loop user study.}
A study with ICS operators would measure whether the V1--V5 severity
classification changes operator behaviour relative to the unguarded
baseline. The study is essential before the framework is deployed in
production and is the most important methodological follow-up.

\paragraph{Multi-protocol generalisation.}
The framework is protocol-agnostic by construction. Generalisation to
DNP3, Modbus, and S7 would broaden the operational applicability and is
the natural next empirical contribution; the bound schema and the
verification pipeline are already protocol-neutral, so the work is
concentrated in the data preprocessing and the event-type classification.
